\item \defpoints{15} [Performing K-Means by Hand]

Let's do $k$-means! To initialize the centroids, we use the $k$-means++ algorithm. And then use Euclidean distance to cluster the 8 data points into $k=3$ clusters. The coordinates of the data points are:
\begin{align*}
    x^{(1)} & = (2,8),  \ x^{(2)} = (2,5), \ x^{(3)} = (1,2), \ x^{(4)} = (5,8), \\
    x^{(5)} & = (7,3),  \ x^{(6)} = (6,4), \ x^{(7)} = (8,4), \ x^{(8)} = (4,7).
\end{align*}
Suppose that initially the first cluster centers is $x^{(1)}$. \\
{\color{blue} To ensure consistent results, please use random numbers in the order shown in the table below. When selecting a center, arrange it in ascending order of sequence number. For example, when the normalized weights of $5$ nodes are $0.2$, $0.1$, $0.3$, $0.3$, and $0.1$, if the random number is $0.3$, the selected node is the third one. Note that you don't necessarily need to use all of them.}
\begin{table*}[h]
    \centering
    \begin{tabular}{|c|c|c|c|c|}
    \hline
    0.6 & 0.2 & 0.5 & 0.9 & 0.3 \\
    \hline
    \end{tabular}
\end{table*}

\begin{itemize}
\item[(a)] Perform the $k$-means++ algorithm to initialize other centers and report the coordinates of the resulting centroids. ~\defpoints{5}
\item[(b)] Calculate the loss function
$$Q(r,c) = \dfrac{1}{n} \sum_{i=1}^n \sum_{j=1}^K r_{ij}\left\|x^{(i)} - c_j\right\|_2^2$$
where $r_{ij} = 1$ if $x^{(i)}$ belongs to the $j$-th cluster and 0 otherwise. ~\defpoints{5}
\item[(c)] How many more iterations after initialization are needed to converge? ~\defpoints{3} Calculate the loss after it converged. ~\defpoints{2}
\end{itemize}


\solution

\textbf{Part (a): K-means++ Initialization}

Given data points:
$x^{(1)} = (2,8)$, $x^{(2)} = (2,5)$, $x^{(3)} = (1,2)$, $x^{(4)} = (5,8)$,
$x^{(5)} = (7,3)$, $x^{(6)} = (6,4)$, $x^{(7)} = (8,4)$, $x^{(8)} = (4,7)$

\textbf{Step 1:} First center is $c_1 = x^{(1)} = (2,8)$.

\textbf{Step 2:} Calculate squared distances from each point to $c_1$:
\begin{align*}
d^2(x^{(1)}, c_1) &= 0\\
d^2(x^{(2)}, c_1) &= (2-2)^2 + (5-8)^2 = 9\\
d^2(x^{(3)}, c_1) &= (1-2)^2 + (2-8)^2 = 1 + 36 = 37\\
d^2(x^{(4)}, c_1) &= (5-2)^2 + (8-8)^2 = 9\\
d^2(x^{(5)}, c_1) &= (7-2)^2 + (3-8)^2 = 25 + 25 = 50\\
d^2(x^{(6)}, c_1) &= (6-2)^2 + (4-8)^2 = 16 + 16 = 32\\
d^2(x^{(7)}, c_1) &= (8-2)^2 + (4-8)^2 = 36 + 16 = 52\\
d^2(x^{(8)}, c_1) &= (4-2)^2 + (7-8)^2 = 4 + 1 = 5
\end{align*}

Sum of distances: $0 + 9 + 37 + 9 + 50 + 32 + 52 + 5 = 194$

Normalized probabilities:
$p_1 = 0/194 = 0$, $p_2 = 9/194 \approx 0.046$, $p_3 = 37/194 \approx 0.191$, $p_4 = 9/194 \approx 0.046$,
$p_5 = 50/194 \approx 0.258$, $p_6 = 32/194 \approx 0.165$, $p_7 = 52/194 \approx 0.268$, $p_8 = 5/194 \approx 0.026$

Cumulative probabilities: $0, 0.046, 0.237, 0.283, 0.541, 0.706, 0.974, 1.0$

Using random number 0.6: This falls in the range for $x^{(6)}$, so $c_2 = (6,4)$.

\textbf{Step 3:} Calculate distances to nearest center for each point:
\begin{align*}
d^2(x^{(1)}) &= \min(0, (2-6)^2 + (8-4)^2) = \min(0, 16 + 16) = 0\\
d^2(x^{(2)}) &= \min(9, (2-6)^2 + (5-4)^2) = \min(9, 16 + 1) = 9\\
d^2(x^{(3)}) &= \min(37, (1-6)^2 + (2-4)^2) = \min(37, 25 + 4) = 29\\
d^2(x^{(4)}) &= \min(9, (5-6)^2 + (8-4)^2) = \min(9, 1 + 16) = 9\\
d^2(x^{(5)}) &= \min(50, (7-6)^2 + (3-4)^2) = \min(50, 1 + 1) = 2\\
d^2(x^{(6)}) &= \min(32, 0) = 0\\
d^2(x^{(7)}) &= \min(52, (8-6)^2 + (4-4)^2) = \min(52, 4 + 0) = 4\\
d^2(x^{(8)}) &= \min(5, (4-6)^2 + (7-4)^2) = \min(5, 4 + 9) = 5
\end{align*}

Sum of distances: $0 + 9 + 29 + 9 + 2 + 0 + 4 + 5 = 58$
Normalized probabilities:
$p_1 = 0/58 = 0$, $p_2 = 9/58 \approx 0.155$, $p_3 = 29/58 \approx 0.500$, $p_4 = 9/58 \approx 0.155$,
$p_5 = 2/58 \approx 0.034$, $p_6 = 0/58 = 0$, $p_7 = 4/58 \approx 0.069$, $p_8 = 5/58 \approx 0.086$

Cumulative probabilities: $0, 0.155, 0.655, 0.810, 0.844, 0.844, 0.913, 1.0$

Using random number 0.2: This falls in the range for $x^{(3)}$, so $c_3 = (1,2)$.

So the final centroids after initialization are:
\begin{align*}
c_1 &= (2,8) \\
c_2 &= (6,4) \\
c_3 &= (1,2)
\end{align*}



\textbf{Part (b): Calculate Initial Loss}

The centroids are:
$c_1 = (2,8)$, $c_2 = (6,4)$, $c_3 = (1,2)$.

The loss function is given by:
$$Q(r,c) = \dfrac{1}{n} \sum_{i=1}^n \sum_{j=1}^K r_{ij}\left\|x^{(i)} - c_j\right\|_2^2$$

Where $r_{ij} = 1$ if $x^{(i)}$ belongs to the $j$-th cluster and 0 otherwise.

We need to assign each point to the nearest centroid and calculate the loss.
\begin{align*}
r_{11} &= 1, \ r_{12} = 0, \ r_{13} = 0 \quad (x^{(1)} \text{ to } c_1)\\
r_{21} &= 1, \ r_{22} = 0, \ r_{23} = 0 \quad (x^{(2)} \text{ to } c_1)\\
r_{31} &= 0, \ r_{32} = 0, \ r_{33} = 1 \quad (x^{(3)} \text{ to } c_3)\\
r_{41} &= 1, \ r_{42} = 0, \ r_{43} = 0 \quad (x^{(4)} \text{ to } c_1)\\
r_{51} &= 0, \ r_{52} = 1, \ r_{53} = 0 \quad (x^{(5)} \text{ to } c_2)\\
r_{61} &= 0, \ r_{62} = 1, \ r_{63} = 0 \quad (x^{(6)} \text{ to } c_2)\\
r_{71} &= 0, \ r_{72} = 1, \ r_{73} = 0 \quad (x^{(7)} \text{ to } c_2)\\
r_{81} &= 1, \ r_{82} = 0, \ r_{83} = 0 \quad (x^{(8)} \text{ to } c_1)
\end{align*}
Calculating the distances:
\begin{align*}
\left\|x^{(1)} - c_1\right\|_2^2 &= 0 \\
\left\|x^{(2)} - c_1\right\|_2^2 &= 9 \\
\left\|x^{(3)} - c_3\right\|_2^2 &= 0 \\
\left\|x^{(4)} - c_1\right\|_2^2 &= 9 \\
\left\|x^{(5)} - c_2\right\|_2^2 &= 2 \\
\left\|x^{(6)} - c_2\right\|_2^2 &= 0 \\
\left\|x^{(7)} - c_2\right\|_2^2 &= 4 \\
\left\|x^{(8)} - c_1\right\|_2^2 &= 5
\end{align*}
Now we can calculate the loss:
\begin{align*}
Q(r,c) &= \frac{1}{8} \left( 0 + 9 + 0 + 9 + 2 + 0 + 4 + 5 \right) \\
&= \frac{29}{8} = 3.625
\end{align*}



\textbf{Part (c): Iterations to Convergence}

\textbf{Iteration 1:}
Update centroids based on current assignments:
\begin{align*}
c_1 &= \left( \frac{2+2+5+4}{4}, \frac{8+5+8+7}{4} \right) = (\frac{13}{4}, 7) \\
c_2 &= \left( \frac{6+7+8}{3}, \frac{4+3+4}{3} \right) = (7, \frac{11}{3}) \\
c_3 &= \left( \frac{1}{1}, \frac{2}{1} \right) = (1, 2)
\end{align*}
\begin{align*}
\text{New centroids: } c_1 = (\frac{13}{4}, 7), \ c_2 = (7, \frac{11}{3}), \ c_3 = (1, 2)
\end{align*}

\begin{align*}
r_{11} &= 1, \ r_{12} = 0, \ r_{13} = 0 \quad (x^{(1)} \text{ to } c_1)\\
r_{21} &= 1, \ r_{22} = 0, \ r_{23} = 0 \quad (x^{(2)} \text{ to } c_1)\\
r_{31} &= 0, \ r_{32} = 0, \ r_{33} = 1 \quad (x^{(3)} \text{ to } c_3)\\
r_{41} &= 1, \ r_{42} = 0, \ r_{43} = 0 \quad (x^{(4)} \text{ to } c_1)\\
r_{51} &= 0, \ r_{52} = 1, \ r_{53} = 0 \quad (x^{(5)} \text{ to } c_2)\\
r_{61} &= 0, \ r_{62} = 1, \ r_{63} = 0 \quad (x^{(6)} \text{ to } c_2)\\
r_{71} &= 0, \ r_{72} = 1, \ r_{73} = 0 \quad (x^{(7)} \text{ to } c_2)\\
r_{81} &= 1, \ r_{82} = 0, \ r_{83} = 0 \quad (x^{(8)} \text{ to } c_1)
\end{align*}
This shows that 1 iteration can already converge, as the assignments remain the same.\\
Calculating the distances:
\begin{align*}
\left\|x^{(1)} - c_1\right\|_2^2 &= \left(2 - \frac{13}{4}\right)^2 + (8 - 7)^2 = \left(\frac{-5}{4}\right)^2 + 1 = \frac{25}{16} + 1 = \frac{41}{16} \\
\left\|x^{(2)} - c_1\right\|_2^2 &= \left(2 - \frac{13}{4}\right)^2 + (5 - 7)^2 = \left(\frac{-5}{4}\right)^2 + 4 = \frac{25}{16} + 4 = \frac{89}{16} \\
\left\|x^{(3)} - c_3\right\|_2^2 &= (1 - 1)^2 + (2 - 2)^2 = 0 \\
\left\|x^{(4)} - c_1\right\|_2^2 &= \left(5 - \frac{13}{4}\right)^2 + (8 - 7)^2 = \left(\frac{7}{4}\right)^2 + 1 = \frac{49}{16} + 1 = \frac{65}{16} \\
\left\|x^{(5)} - c_2\right\|_2^2 &= (7 - 7)^2 + \left(3 - \frac{11}{3}\right)^2 = 0 + \left(\frac{-2}{3}\right)^2 = \frac{4}{9} \\
\left\|x^{(6)} - c_2\right\|_2^2 &= (6 - 7)^2 + \left(4 - \frac{11}{3}\right)^2 = 1 + \left(\frac{1}{3}\right)^2 = 1 + \frac{1}{9} = \frac{10}{9} \\
\left\|x^{(7)} - c_2\right\|_2^2 &= (8 - 7)^2 + \left(4 - \frac{11}{3}\right)^2 = 1 + \left(\frac{1}{3}\right)^2 = 1 + \frac{1}{9} = \frac{10}{9} \\
\left\|x^{(8)} - c_1\right\|_2^2 &= \left(4 - \frac{13}{4}\right)^2 + (7 - 7)^2 = \left(\frac{3}{4}\right)^2 + 0 = \frac{9}{16}
\end{align*}
Calculating the new loss:
$Q(r,c) = \frac{1}{8} \left( \frac{41}{16} + \frac{89}{16} + 0 + \frac{65}{16} + \frac{4}{9} + \frac{10}{9} + \frac{10}{9} + \frac{9}{16} \right) = \frac{185}{96}$






\newpage